\section{検証データと目視ラベルの再現性}
\label{sec_validation}
\subsection{検証データの作成}
検証データは，定義書の改訂に合わせて2回に分けて作成した．
最初に，\ref{sec_label}章の7つの層から各12件のお題を抽出し，各お題から評価の上位3件と無作為の2件のボケを取って420件とした．
この420件では，セリフと回答が大半を占め，ツッコミは12件，心の声は28件しかなかった．
そこで，教師または規則が心の声，ツッコミ，状況の説明，名付け・見立て，セリフのいずれかと予測したボケから，型ごとに40件を予測を見せずに追加し，計620件とした．
ツッコミの候補は，教師も規則もツッコミと予測したボケであったが，目視では多くがセリフか状況の説明であり，追加後もツッコミは26件にとどまった．

620件は，お題の単位で層ごとに半数ずつ，調整用データ（dev，308件）と最終評価用データ（test，312件）に分けた．
指示文や構成の選択には調整用データだけを用い，最終評価用データは精度の報告にのみ用いた．
評価指標には，各クラスの F1 の平均であるマクロ F1 を用いた．
クラス $ k $ の適合率を $ P_k $，再現率を $ R_k $ とすると，マクロ F1 は(\ref{eq_macro_f1})で表される．

\begin{equation}
    {\rm F1}_{\rm macro} = \frac{1}{K} \sum_{k=1}^{K} \frac{2 P_k R_k}{P_k + R_k}
    \label{eq_macro_f1}
\end{equation}

\noindent
ここで $ K $ はクラス数である．
画像外の知識は2値であるため，「はい」のクラスの F1 を用いた．
最終評価用データの F1 の95\%信頼区間は，項目を復元抽出する1,000回のブートストラップで求めた．

\subsection{目視ラベルの見直し}
目視ラベルは同一の判断者によるものであり，判断の揺れや系統的な誤りを含みうる．
最初の420件では，2つの教師がともに目視と異なった項目を定義書に照らして見直した．
対象は必要な知識143件，発話の型109件，不適切19件であったが，大半は教師の系統的な偏りであり，定義書が教師を支持した38件だけを改めた．
第3版で必要な知識を2値に改めた際には，620件の全てを新しい定義で見直し，27件を改めた．
画像外の知識が「はい」の項目は620件中139件であった．

\subsection{付け直しによる再現性の測定}
自動判定の精度がどこまで上がりうるかを知るため，620件の全てを無作為に並べ替え，元の番号とラベルを伏せた一覧で，定義書 第3版に従って付け直した．
付け直しは，同じ判断者が時間を置いて行ったものであり，判断者の間の一致度ではなく，同一の判断者の再現性を表す．
表\ref{tab_reliability}に，元のラベルとの一致を示す．

\begin{table}[t]
    \centering
    \caption{元の目視ラベルと付け直したラベルの一致（620件）．$\kappa$ は Cohen の $\kappa$ 係数\cite{ref_kappa}である．}
    \label{tab_reliability}
    \small
    \begin{tabular}{lccc}
        \toprule
        項目 & F1 & 一致率 & $\kappa$ \\
        \midrule
        発話の型 6区分（マクロ） & 0.80 & 0.86 & 0.82 \\
        発話の型 3区分（マクロ） & 0.93 & --- & --- \\
        画像外の知識 & 0.85 & --- & 0.81 \\
        \bottomrule
    \end{tabular}
\end{table}

6区分のクラスごとの F1 は，回答が0.94，セリフが0.89，状況の説明が0.85，名付け・見立てが0.81である一方，心の声は0.75，ツッコミは0.55であった．
不一致で多かったのは，心の声とセリフ（11件），セリフとツッコミ，状況の説明と名付け・見立ての間の揺れであった．
また，形式が通常でも，お題画像に空欄や空の吹き出し，問いが含まれる場合に回答とするかの揺れが約10件あった．

不一致の項目は定義書に照らして裁定し，合意ラベルとした．
裁定では，お題画像が空欄・空の吹き出し・問いを課している場合は形式によらず回答とし，迷う場合はセリフとし，画像外の知識で判断が分かれた項目は元のラベルを維持した．
合意ラベルは元のラベルから，発話の型で55件，画像外の知識で22件が変わった．
以下の評価は，特に断らない限り，全ての方法をこの合意ラベルで評価し直した値である．
そのため，各実験の当時に報告した値とは小数第2位で異なる場合がある．

表\ref{tab_reliability}の一致度は，同じ定義でも時間を置くと6区分の約14\%の判断が変わり，少ないクラスほど揺れが大きいことを示す．
したがって，自動判定の F1 が6区分で0.8，画像外の知識で0.85を大きく超えることは期待しにくく，これらの値を実質的な上限とみなす．
